\section{Лемма Цорна}

\begin{definition}[6.1]
\begin{enumerate}[label=(\alph*)]
\item Отношение $\le$ на множестве $S$~--- \emph{частичный порядок}, если оно рефлексивно ($a\le a$), транзитивно
($a\le b$ и $b\le c$ $\Rightarrow$ $a\le c$) и антисимметрично ($a\le b$ и $b\le a$ $\Rightarrow$ $a=b$).
\item Частичный порядок \emph{линеен}, если любые два элемента сравнимы: $s\le t$ или $t\le s$. Подмножество
частично упорядоченного множества, линейно упорядоченное индуцированным порядком, называется \emph{цепью}.
\item \emph{Верхняя грань} подмножества $T$ частично упорядоченного множества $(S,\le)$~--- такой элемент $s\in S$,
что $t\le s$ для всех $t\in T$.
\item \emph{Максимальный элемент} частично упорядоченного множества $S$~--- такой элемент $s\in S$, что
$s\le s'$ $\Rightarrow$ $s=s'$.
\end{enumerate}
\end{definition}

Максимальных элементов может не быть: множество конечных подмножеств бесконечного множества, упорядоченное по
включению, их не имеет. Максимальный элемент не обязан быть наибольшим: в множестве собственных подпространств
плоскости $\R^2$, упорядоченном по включению, все прямые максимальны.

\begin{lemma}[6.2: лемма Цорна]
Пусть $(S,\le)$~--- непустое частично упорядоченное множество, в котором у каждой цепи есть верхняя грань. Тогда в $S$
есть максимальный элемент.
\end{lemma}

\emph{Аксиома выбора} утверждает, что для любого семейства $(X_i)_{i\in I}$ непустых множеств существует
\emph{функция выбора}~--- отображение $i\mapsto x_i\in X_i$; иначе говоря, произведение $\prod_iX_i$ непусто. Лемма
Цорна равносильна аксиоме выбора и потому не зависит от аксиом теории множеств Цермело--Френкеля (ZF); мы принимаем её
как аксиому. Следующее предложение~--- типичное применение.

\begin{proposition}[6.3]
В каждом ненулевом коммутативном кольце $A$ есть максимальный идеал (максимальный среди собственных идеалов).
\end{proposition}

\begin{proof}
Пусть $S$~--- множество собственных идеалов $A$, упорядоченное по включению; оно непусто: идеал $0$ собственный, так как
$A\ne0$. Пустая цепь ограничена идеалом $0$; пусть $T$~--- непустая цепь идеалов. Тогда $J=\bigcup_{I\in T}I$~--- снова идеал: любые два элемента $J$ лежат в одном идеале из $T$
(цепь линейно упорядочена), а он замкнут относительно сложения и умножения на элементы $A$. Идеал $J$ собственный:
если $1\in J$, то $1\in I$ для некоторого $I\in T$, а $I$ собственный. Значит, $J$~--- верхняя грань $T$ в $S$. По лемме Цорна в $S$ есть максимальный элемент; это и есть
максимальный идеал.
\end{proof}

\begin{remark}[6.4]
Лемма Цорна на самом деле равносильна существованию максимальных идеалов в ненулевых коммутативных кольцах. Более
слабая аксиома~--- \emph{принцип ультрафильтра} (каждый собственный фильтр на множестве содержится в
ультрафильтре)~--- равносильна
существованию простых идеалов в ненулевых коммутативных кольцах и компактности произведений компактных хаусдорфовых
пространств; из неё следует аксиома выбора для семейств конечных множеств. Её достаточно для основных результатов
главы~--- существования и единственности алгебраического замыкания (первое и третье доказательства, \rf{6.7}); см.~\rf{6.10a}.
\end{remark}

\begin{aside}[6.4a: формы аксиомы выбора]
Вот как связаны утверждения, которые встречаются в этой и следующей главах.
\begin{itemize}[leftmargin=*]
\item \emph{Равносильны аксиоме выбора} (при остальных аксиомах ZF): лемма Цорна; теорема Цермело (каждое множество
можно вполне упорядочить); сравнимость любых двух мощностей (ср.~\rf{6.4e}); равенство $|X\times X|=|X|$ для всех
бесконечных $X$ (ср.~\rf{6.4f}); существование максимальных идеалов (ср.~\rf{6.3}); существование базиса у каждого векторного
пространства; теорема Тихонова о компактности произведения произвольных (не обязательно хаусдорфовых) компактных
пространств.
\item \emph{Строго слабее}~--- принцип ультрафильтра. Ему равносильны существование простых идеалов в ненулевых
коммутативных кольцах и компактность произведения компактных хаусдорфовых пространств. Из него следуют аксиома выбора
для семейств конечных множеств, существование и единственность (с точностью до изоморфизма) алгебраического замыкания
(первое и третье доказательства ниже, \rf{6.7}) и компактность групп Галуа (\rf{7.8}).
\item \emph{Без выбора} этих результатов не доказать. С аксиомами ZF совместимо существование алгебраического замыкания
поля $\Q$, у которого нет нетривиальных автоморфизмов; оно не изоморфно алгебраическому замыканию $\Q$ в $\CC$, где есть
комплексное сопряжение. Значит, в ZF единственность алгебраического замыкания недоказуема. Недоказуемо и
существование: с ZF совместимо, что у некоторого поля нет алгебраического замыкания.
\end{itemize}
\end{aside}
